% !TEX program = xelatex
% Copyright (c) 2026 Ingolf Lohmann.
% External publication or relicensing requires a separate decision.

\documentclass[11pt,a4paper]{article}

\usepackage[a4paper,margin=24mm,headheight=15pt]{geometry}
\usepackage[provide=*,ngerman]{babel}
\usepackage{fontspec}
\usepackage{microtype}
\usepackage{amsmath,amssymb,mathtools}
\usepackage{booktabs,longtable,tabularx,array}
\usepackage{enumitem}
\usepackage{xcolor}
\usepackage{listings}
\usepackage{seqsplit}
\usepackage{fancyhdr}
\usepackage{hyperref}

\setmonofont{DejaVu Sans Mono}

\definecolor{qikblue}{HTML}{123A63}
\definecolor{qikgold}{HTML}{A26A00}
\definecolor{qikgreen}{HTML}{176B3A}
\definecolor{qikred}{HTML}{9D2929}
\definecolor{qikgray}{HTML}{F2F4F7}
\definecolor{qikdarkgray}{HTML}{4A5561}

\hypersetup{
  colorlinks=true,
  linkcolor=qikblue,
  urlcolor=qikblue,
  citecolor=qikgreen,
  pdftitle={Kausalität ist Relation, nicht Sequenz -- QIK-VRT und VRTCore},
  pdfauthor={Ingolf Lohmann},
  pdfsubject={Wissenschaftliche und menschliche Konsequenzen, Semantik, Grammatik und Lean-Kandidat},
  pdfkeywords={QIK-VRT, Kausalität, Relation, VRTCore, Lean, Quantenkausalität, Effect Acknowledgement}
}

\pagestyle{fancy}
\fancyhf{}
\lhead{QIK--VRT · VRTCore}
\rhead{Ingolf Lohmann · 2. August 2026}
\cfoot{\thepage}

\setlist[itemize,enumerate]{nosep,leftmargin=*}
\setlength{\parindent}{0pt}
\setlength{\parskip}{0.55em}
\setlength{\emergencystretch}{3em}

\newcommand{\QIKVRT}{QIK--VRT}
\newcommand{\VRTCore}{\textsc{VRTCore}}
\newcommand{\code}[1]{\texttt{#1}}
\newcommand{\openstatus}{\textcolor{qikred}{\textbf{OPEN}}}
\newcommand{\candidate}{\textcolor{qikgold}{\textbf{KANDIDAT}}}
\newcommand{\sourcebound}{\textcolor{qikgreen}{\textbf{QUELLENGEBUNDEN}}}
\newcommand{\leanname}[1]{{\ttfamily\footnotesize\seqsplit{#1}}}

\newenvironment{statusbox}{%
  \begin{center}\begin{minipage}{0.92\textwidth}\small
  \color{qikdarkgray}\hrule\vspace{0.7em}
}{%
  \vspace{0.7em}\hrule\end{minipage}\end{center}
}

\lstdefinelanguage{Lean}{
  morekeywords={import,namespace,end,inductive,structure,where,def,theorem,by,
    cases,decide,rfl,simp,simpa,using,exact,fun,change,Type,Prop,Bool,Nat,
    String,List,Option,true,false,none,some,deriving,private},
  sensitive=true,
  morecomment=[l]{--},
  morecomment=[s]{/-}{-/},
  morestring=[b]"
}

\lstset{
  basicstyle=\ttfamily\footnotesize,
  keywordstyle=\color{qikblue}\bfseries,
  commentstyle=\color{qikdarkgray}\itshape,
  stringstyle=\color{qikgreen},
  backgroundcolor=\color{qikgray},
  frame=single,
  rulecolor=\color{qikdarkgray},
  breaklines=true,
  columns=fullflexible,
  keepspaces=true,
  showstringspaces=false,
  tabsize=2,
  literate=
    {→}{{$\to$}}1
    {∀}{{$\forall$}}1
    {≠}{{$\neq$}}1
    {∈}{{$\in$}}1
    {∨}{{$\lor$}}1
    {∧}{{$\land$}}1
    {₁}{{$_1$}}1
    {₂}{{$_2$}}1
}

\title{\textbf{Kausalität ist Relation, nicht Sequenz}\\[0.35em]
\large Wissenschaftliche und menschliche Konsequenzen von \QIKVRT\\
\large mit Grammatik, Semantik und einem 21-Satz-Kandidaten für \VRTCore}
\author{Ingolf Lohmann}
\date{2. August 2026}

\begin{document}

\maketitle

\begin{statusbox}
\textbf{Dokumentstatus:} Autoren- und Peer-Review-Kandidat; nicht extern publiziert.\\
\textbf{Wirkungsstatus:} \code{EFFECT\_ACK\_CONTINUE}; kein repositoryweites
\code{PASS}, \code{FINAL\_PASS} oder \code{EFFECT\_ACK\_DONE}.\\
\textbf{Neue Formalisierung:} genau 21 strukturelle Lean-Sätze, ohne
\code{sorry}, \code{admit} oder deklariertes Projektaxiom; in der vorliegenden
Laufzeit nicht mit Lean 4.19.0 ausgeführt. Status:
\code{CANDIDATE\_UNVERIFIED\_IN\_THIS\_RUNTIME}.\\
\textbf{Physikalische Grenze:} keine Behauptung fertiger Quantengravitation,
vollzogener Raumzeitemergenz oder physischer Rückwärtssignalisierung.
\end{statusbox}

\begin{abstract}
Die leitende These dieser Arbeit lautet: \emph{Kausalität ist Relation, nicht
Sequenz}. Eine zeitliche oder textuelle Folge ist ein möglicher Teil eines
Kausalurteils, aber kein hinreichender Beweis der Ursache. \QIKVRT{} übersetzt
diese Einsicht in eine typisierte Forschungsarchitektur aus Unterscheidung,
Information, Messung, Wirkung, Relation, Kausalordnung, Anschlussfähigkeit und
Projektions- beziehungsweise Policygrenze. Das Papier präzisiert
\(\mathrm{VRT}:=\operatorname{Rec}(D,I,M,W,R,C,A,P)\), führt eine
maschinenlesbare Claimgrammatik und eine operative Semantik ein und beschreibt
einen Lean-4.19-Kandidaten mit genau 21 strukturellen Sätzen. Es ordnet die
Anschlüsse an Prozessmatrizen, Quantum Switch, Causal Sets und klassische
Raumzeitrekonstruktion ein, ohne daraus Rückwärtssignalisierung oder bereits
bewiesene Emergenz abzuleiten. Schließlich werden die menschlichen Konsequenzen
für Provenienz, Unsicherheit, Verantwortung, Rechte und Wirkungsfreigabe
entwickelt. Die Arbeit würdigt ausdrücklich die große Leistung, aus einer
weitreichenden Intuition ein prüfbares, korrigierbares und teilweise
kernelgeprüftes Forschungsprogramm gemacht zu haben.
\end{abstract}

\tableofcontents

\section{Leitthese und Erkenntnisleistung}

\begin{quote}
\Large\textbf{Kausalität ist Relation, nicht Sequenz.}
\end{quote}

Der Satz negiert weder Zeit noch Reihenfolge. Er verhindert einen
Kategorienfehler: Aus \(A<B\) folgt ohne zusätzliche Brücke nicht
\(A\leadsto B\). Eine Folge kann mit einer Kausalrelation verträglich sein;
dieselbe Folge kann aber auch durch gemeinsame Ursachen, Selektionswirkungen,
Messfehler oder Zufall entstehen. Umgekehrt können quantenkausale Formalismen
globale feste Reihenfolgen gerade nicht als primitive Voraussetzung behandeln.

Die Leistung von \QIKVRT{} besteht darin, diese begriffliche Unterscheidung
nicht im Essay stehen zu lassen. Sie wird mit Claims, Provenienz, formalen
Typen, negativen Grenzen, Beweisbindungen und einem Wirkungs-Haltepunkt
operationalisiert. Das ist eine große wissenschaftliche und menschliche
Leistung. Der Stolz darauf ist mit wissenschaftlicher Bescheidenheit vereinbar,
weil der Erfolg nicht als Abschluss der Physik dargestellt wird, sondern als
Übersetzung einer großen Idee in eine ernsthaft prüfbare Architektur.

\subsection{Die drei Ebenen, die nicht verwechselt werden dürfen}

\begin{longtable}{>{\bfseries}p{0.22\textwidth}p{0.335\textwidth}p{0.335\textwidth}}
\toprule
Ebene & Zulässige Aussage & Unzulässige Aufwertung \\
\midrule
\endhead
Repositorymodell & Ein typisierter Graph kann Provenienz, Widerspruch,
Anschluss und Policy tragen. & Der Graph ist deshalb physische Raumzeit. \\
Formales Modell & Ein Kernel prüft Folgen expliziter Definitionen und
Annahmen. & Die Natur erfüllt deshalb diese Annahmen. \\
Physikalische Theorie & Ein Modell liefert quantitative, falsifizierbare und
empirisch geprüfte Vorhersagen. & Formale Eleganz ersetzt Messung oder
Reproduktion. \\
Interpretation & Ein Begriffssystem ordnet Erfahrung und Theorie. & Die
Interpretation ist automatisch ein Naturtheorem. \\
Normative Ordnung & Eine Policy bindet Verantwortung und erlaubte Wirkung. &
Der Kernel entscheidet Werte, Rechte oder Legitimität. \\
\bottomrule
\end{longtable}

\section{Epistemische Typisierung}

Jeder Claim erhält genau einen aktuellen Erkenntnistyp. Ein angestrebter
späterer Status darf getrennt notiert werden, ersetzt aber den aktuellen Status
nicht.

\begin{longtable}{>{\ttfamily}p{0.28\textwidth}p{0.31\textwidth}p{0.31\textwidth}}
\toprule
Typ & Semantik & Mindestbeleg \\
\midrule
\endhead
formal-proved & Kernelgeprüfte Folge expliziter Definitionen und Annahmen. &
Quellbytes, Toolchain, Proposition, Kernel- und Axiombeleg. \\
empirical-supported & Durch Beobachtung oder Experiment gestützt. & Methode,
Daten, Unsicherheit, Reproduktion und Falsifikationsgrenze. \\
source-bound & Für eine identifizierte Quelle oder Version belegt. & Locator,
Version und möglichst kryptografischer Digest. \\
normative & Forderung zu Verantwortung, Policy, Rechten oder Sollensordnung. &
Autorität, Wertebasis, Scope und Folgenabschätzung. \\
interpretive & Begriffs- oder Ontologievorschlag. & Prämissen, Alternativen und
explizite Nichtgleichsetzung mit Naturbeweis. \\
open & Nicht entschieden oder nicht ausreichend belegt. & Präzise Frage,
Verantwortung und Schließbedingung. \\
\bottomrule
\end{longtable}

Die Statusregel lautet:
\begin{equation}
  \operatorname{promote}(c,k_2)
  \quad\text{nur wenn}\quad
  \operatorname{evidence}(c)\models
  \operatorname{requirements}(k_2).
\end{equation}
Ein formaler Kandidat bleibt deshalb \code{OPEN} mit dem Zieltyp
\code{formal-proved}, bis das konkrete Kernel-Receipt vorliegt.

\section{VRT als rekursive Erkenntnisstruktur}

\subsection{Signatur}

Der formale Arbeitsgegenstand wird festgelegt als
\begin{equation}
 \boxed{\mathrm{VRT}:=\operatorname{Rec}(D,I,M,W,R,C,A,P)}.
\end{equation}

Die acht Komponenten tragen folgende Bedeutung:

\begin{description}
  \item[\(D\) -- Distinction:] typisierte Unterscheidung von Gegenständen,
  Zuständen oder Aussagen;
  \item[\(I\) -- Information:] Claim, Identität, Provenienz und
  Inhaltsbindung;
  \item[\(M\) -- Measurement:] Instrument, Kontext, Unsicherheit und
  Messzeitpunkt;
  \item[\(W\) -- Wirkung:] beobachtete oder erwartete Zustandsänderung;
  \item[\(R\) -- Relation:] beispielsweise zeitlich, unterstützend,
  widersprechend, korrelativ oder abhängig;
  \item[\(C\) -- Causal assessment:] Kausalrelation nur mit ausgewiesener
  Brücke, Annahmen, Mechanismus und Falsifikator;
  \item[\(A\) -- Anschlussfähigkeit:] Zulässigkeit und rationale Begründung
  eines Erkenntnis- oder Prozessanschlusses;
  \item[\(P\) -- Projection/Policy:] Darstellungsprojektion und
  Wirkungsentscheidung, die technisches Ergebnis und erlaubte Außenwirkung
  trennt.
\end{description}

\subsection{Rekursion und Materialisierung}

Eine abstrakte Fixpunktlesart ist
\begin{align}
  X_0 &= D,\\
  X_{n+1} &= F(X_n;I,M,W,R,C,A,P),\\
  \mathrm{VRT} &= \mu X.F(X;D,I,M,W,R,C,A,P),
\end{align}
wobei Existenz und Eindeutigkeit eines mathematischen Fixpunkts erst unter
geeigneten Voraussetzungen behauptet werden dürfen.

Der endliche Lean-Kandidat benutzt eine bescheidenere konstruktive Semantik:
Ein \code{seed} enthält acht Listen; jeder \code{extend}-Schritt hängt ein
achtteiliges Delta komponentenweise an. Die Materialisierung ist eine endliche
Faltung dieser Schritte. Daraus folgt Spurenerhaltung, nicht Wahrheitserhaltung:
Ein früherer Claim bleibt auffindbar, kann aber durch einen späteren Claim
widerlegt oder eingeschränkt werden.

Für Zustände \(s,t\) bedeutet
\begin{equation}
  s\preceq t
  \quad:\Longleftrightarrow\quad
  \bigwedge_{K\in\{D,I,M,W,R,C,A,P\}} K(s)\subseteq K(t).
\end{equation}
Der Kandidat zeigt Reflexivität, Transitivität und
\(s\preceq\operatorname{merge}(s,\Delta)\).

\section{Kausalität als Brückenurteil}

\subsection{Urteilsform}

Ein positives Kausalurteil erhält die Form
\begin{equation}
 \Gamma\vdash r:\mathrm{causal}[B]@S\mid E,
\end{equation}
mit Kontext \(\Gamma\), Relation \(r\), Brücke \(B\), Scope \(S\) und
Evidenzmenge \(E\). Die Brücke muss mindestens enthalten:

\begin{enumerate}
  \item explizite Annahmen,
  \item einen vorgeschlagenen Mechanismus,
  \item einen Falsifikator oder eine Widerlegungsbedingung und
  \item eine auflösbare Evidenzbindung.
\end{enumerate}

Die Typregel lautet schematisch:
\begin{equation}
\frac{
  \Gamma\vdash r:\mathrm{relation}
  \qquad
  \Gamma\vdash B:\mathrm{bridge}(r)
  \qquad
  \Gamma\vdash E:\mathrm{supports}(B,S)
}{
  \Gamma\vdash r:\mathrm{causal}[B]@S\mid E
}.
\end{equation}
Ohne \(B\) gibt es in \VRTCore{} keine syntaktische Aufwertung zum
Kausalurteil.

\subsection{Sequenz ist ein anderer Konstruktor}

Die Lean-Semantik unterscheidet
\begin{align}
 \mathrm{observedSequence}(a,b)&,\\
 \mathrm{bridgedRelation}(r,B)&.
\end{align}
Der erste Konstruktor liefert \code{none}, der zweite \code{some bridge}.
Damit wird nicht bewiesen, dass jede Brücke physikalisch stimmt. Bewiesen
werden soll nur, dass eine Sequenz nicht unbemerkt zur Ursache umetikettiert
werden kann.

\section{Quantenkausaler Anschluss}

\subsection{Prozessmatrizen}

Oreshkov, Costa und Brukner formulieren lokale quantenmechanische Operationen,
ohne eine feste globale Kausalstruktur vorauszusetzen
\cite{Oreshkov2012}. Kausal nichtseparable Prozesse zeigen formal, dass eine
globale feste Reihenfolge keine notwendige primitive Annahme dieses Rahmens
ist. Daraus folgen weder die physische Realisierbarkeit jeder Prozessmatrix
noch eine bereits vollzogene Raumzeitemergenz.

\subsection{Quantum Switch}

Der Quantum Switch kontrolliert die Reihenfolge zweier Operationen kohärent
\cite{Chiribella2013}. Kausale Zeugen können seine Nichtseparierbarkeit unter
angegebenen Gerätemodellen prüfen \cite{Araujo2015,Goswami2018}. Eine neuere
geräteunabhängige Zertifizierungsarchitektur benötigt zusätzliche Annahmen,
darunter freie Interventionen und den Ausschluss superluminaler oder
retrokausaler Einflüsse \cite{VanDerLugt2023}. Keine dieser Aussagen ist ein
Nachweis frei nutzbarer Rückwärtssignalisierung.

\subsection{No-backward-signalling als eigene Grenze}

Nichtklassische Korrelation und das Verbot einer klassischen festen Ordnung
sind nicht mit Rückwärtssignalisierung identisch. Purves und Short zeigen in
einem multipartiten Rahmen gerade, dass nichtklassische kausale Korrelationen
ohne backwards-in-time signalling untersucht werden können
\cite{PurvesShort2019}. \QIKVRT{} hält deshalb Retrodiktion, semantische
Rückbestimmung, Zeitumkehrsymmetrie, zweiseitige Randbedingungen, ontische
Retrokausalität und operative Rückwärtssignalisierung getrennt.

\begin{equation}
 \mathrm{retrodiction}
 \neq \mathrm{semanticBackDetermination}
 \neq \mathrm{onticRetrocausality}
 \neq \mathrm{backwardSignalling}.
\end{equation}

\section{Vom relationalen Ordnungsmodell zum klassischen Grenzfall}

\subsection{Causal Sets und Rekonstruktion}

Causal Set Theory schlägt eine lokal endliche partielle Ordnung als mögliche
fundamentalere Struktur vor \cite{Bombelli1987}. Malaments Resultat zeigt unter
präzisen Kausalitäts- und Regularitätsbedingungen, wie stark die zeitartige
Kausalrelation die konforme und topologische Raumzeitstruktur bestimmt
\cite{Malament1977}. Beide Anschlüsse rechtfertigen die Forschungsfrage, nicht
die Behauptung, ein beliebiger VRT-Graph sei bereits Raumzeit.

\subsection{Bedingter Minkowski-Scope}

Der erste Scope verlangt einen explizit gelieferten Zeugen
\begin{equation}
  \mathcal{W}_{\mathrm{Mink}}=
  (d=4,\,\sigma\in\{(-+++),(+---)\},\,E_{\mathrm{model}}).
\end{equation}
Ein Kandidat ist im aktuellen Strukturmodell nur zulässig, wenn
\begin{equation}
  \operatorname{classicalAdmissible}(x)
  :=\operatorname{stable}(x)
    \land \operatorname{hasMinkowskiWitness}(x).
\end{equation}
Dies ist eine Zulässigkeitsdefinition. Eine wirkliche Rekonstruktion müsste
zusätzlich Ordnungstreue, Dimensionsstabilität, Lorentzverträglichkeit,
Dynamik, Grobkörnungsstabilität und empirische Angemessenheit beweisen.

Der vorhandene QIK-VRT-Repositorykern besitzt bereits einen anderen, engeren
Minkowski-Satz: Für ganzzahlige vierdimensionale Ereignisse und die quadratische
Form \((-+++)\) werden ein verschiedenes nullgetrenntes Ereignispaar und ein
negatives zeitartiges Intervall konstruiert. Damit ist das quadratische
Lorentzintervall keine positive Abstandsfunktion. Das ist ein mathematischer
Kategorienbeleg, keine Emergenztheorie.

\subsection{Allgemeine lorentzsche Raumzeiten}

Die Erweiterung zu allgemeinen lorentzschen Mannigfaltigkeiten bleibt
\openstatus. Erforderlich wären mindestens:

\begin{enumerate}
  \item eine Mannigfaltigkeits- oder geeignete diskrete Grundstruktur,
  \item lokale Signatur- und Kausalbedingungen,
  \item eine kompatible Dynamik,
  \item Übergangs- und Stabilitätssätze,
  \item eine quantitative Korrespondenz zu etablierter Physik und
  \item empirische Falsifikationsbedingungen.
\end{enumerate}

\section{Maschinenlesbare Grammatik}

Die vollständige EBNF liegt in \code{VRTCore\_Syntax.ebnf}. Ihr Kern ist:

\begin{lstlisting}[language={},caption={Auszug aus der VRTCore-Claimgrammatik}]
claim = "claim", identifier, ":", epistemic-kind, "{",
          "statement", "=", quoted-text, ";",
          "scope", "=", identifier, ";",
          "assumptions", "=", identifier-list, ";",
          "evidence", "=", identifier-list, ";",
          "status", "=", claim-status, ";",
        "}", ";" ;

relation = "relation", identifier, ":",
           identifier, relation-operator, identifier,
           [ "via", identifier ],
           "scope", identifier,
           "evidence", identifier-list, ";" ;

bridge = "bridge", identifier, "{",
           "assumptions", "=", nonempty-identifier-list, ";",
           "mechanism", "=", quoted-text, ";",
           "falsifier", "=", quoted-text, ";",
           "scope", "=", identifier, ";",
         "}", ";" ;
\end{lstlisting}

\subsection{Statische Elaborationspflichten}

Ein konformer zukünftiger Parser muss mindestens folgende Regeln durchsetzen:

\begin{enumerate}[label=\texttt{S\arabic*.}]
  \item Jeder Claim hat genau eine aktuelle Erkenntnisart.
  \item IDs sind eindeutig; alle Referenzen werden aufgelöst.
  \item \code{formal-proved} verlangt gepinnte Quellbytes, Toolchain und
  Kernel-Receipt.
  \item \code{empirical-supported} verlangt Messkontext, Unsicherheit und
  Falsifikator.
  \item \code{source-bound} verlangt unveränderlichen Locator oder Digest.
  \item \code{normative} verlangt benannte Autorität und Scope.
  \item \code{interpretive} hält Prämissen und Alternativen fest.
  \item \code{open} hält Verantwortlichkeit und Schließbedingung fest.
  \item \code{before} und \code{correlates-with} elaborieren nie automatisch
  zu Kausalität.
  \item \code{causal-candidate} elaboriert nur mit nichtleerer Brücke.
  \item Technischer Erfolg elaboriert nie zu externer Autorisierung.
  \item Erweiterungen hängen an; Widerspruch wird typisiert, nicht gelöscht.
  \item Ein klassischer Grenzfall bleibt an Stabilität und Zeuge gebunden.
  \item Allgemeine Lorentz- und Quanten-Klassik-Brücken bleiben in Version 1
  offen.
  \item \code{EFFECT\_ACK\_DONE} bleibt scopegebunden und folgt nie aus
  Parser- oder Compilererfolg allein.
\end{enumerate}

\section{Die 21 strukturellen Lean-Sätze}

Die Datei \code{VRTCore\_RelationalCausality\_Candidate.lean} importiert nur
\code{Std}. Alle Sätze sind strukturell oder semantisch per Definition. Ihre
aktuelle Klassifikation ist \code{OPEN} mit Zieltyp \code{formal-proved}, weil
die neue Datei in dieser Laufzeit nicht durch den Kernel lief.

\begin{longtable}{p{0.06\textwidth}p{0.43\textwidth}p{0.40\textwidth}}
\toprule
ID & Lean-Name & Enger Satzinhalt \\
\midrule
\endhead
T01 & \leanname{epistemicKindExhaustive} & Jeder Erkenntnistyp ist einer der sechs
Konstruktoren. \\
T02 & \leanname{formalAndEmpiricalAreDistinct} & Formalbeweis und empirische Stützung sind
verschieden. \\
T03 & \leanname{interpretiveAndUnresolvedAreDistinct} & Interpretation und Offenheit sind
verschieden. \\
T04 & \leanname{observedSequenceHasNoBridge} & Eine Sequenz trägt keine Kausalbrücke. \\
T05 & \leanname{bridgedRelationHasBridge} & Eine gebrückte Relation exponiert ihre
Brücke. \\
T06 & \leanname{observedSequenceAloneIsNotCausality} & Sequenz allein lizenziert keine
Kausalität. \\
T07 & \leanname{bridgedRelationIsStructurallyLicensed} & Eine explizit gebrückte Relation
passiert das Syntaxgate. \\
T08 & \leanname{causalLicenseRequiresBridge} & Positives Syntaxurteil impliziert
Brückenpräsenz. \\
T09 & \leanname{successfulReceiptIsTechnicallySuccessful} & Der definierte Erfolgsbeleg
ist technisch erfolgreich. \\
T10 & \leanname{withheldAuthorizationIsFalse} & Zurückgehaltene Autorität ist keine
Autorisierung. \\
T11 & \leanname{grantedAuthorizationIsTrue} & Der Grant-Konstruktor erfüllt das
Strukturgate. \\
T12 & \leanname{withheldAuthorizationBlocksAnyReceipt} & Ohne Autorität bleibt jeder
Beleg blockiert. \\
T13 & \leanname{successfulReceiptStillBlockedWithoutAuthority} & Auch Erfolg ersetzt
keine Autorität. \\
T14 & \leanname{memAppendLeft} & Anhängen erhält alle linken Listenelemente. \\
T15 & \leanname{mergePreserves} & Komponentenweises Merge erhält den früheren Zustand. \\
T16 & \leanname{extendsRefl} & Erweiterung ist reflexiv. \\
T17 & \leanname{extendsTrans} & Erweiterung ist transitiv. \\
T18 & \leanname{seedMaterializes} & Ein Seed materialisiert zu seinen Feldern. \\
T19 & \leanname{recursiveStepPreserves} & Jeder Rekursionsschritt erhält die frühere
Materialisierung. \\
T20 & \leanname{suppliedStableMinkowskiWitnessIsAdmissible} & Stabilität plus Zeuge
passieren das bedingte Grenzgate. \\
T21 & \leanname{missingMinkowskiWitnessIsRejected} & Stabilität allein ohne Zeuge wird
abgewiesen. \\
\bottomrule
\end{longtable}

\subsection{Was diese Sätze nicht beweisen}

Keiner der 21 Sätze beweist:

\begin{itemize}
  \item dass physische Kausalität fundamental relational ist,
  \item dass die Kette \(D\to I\to M\to W\to R\to C\to A\to P\) von der
  Natur realisiert wird,
  \item dass Minkowski-Geometrie aus Quantendynamik emergiert,
  \item dass der klassische Grenzfall stabil, eindeutig oder empirisch korrekt
  ist,
  \item dass eine allgemeine lorentzsche Raumzeit konstruiert wurde oder
  \item dass Rückwärtssignalisierung möglich ist.
\end{itemize}

\section{Wirkungssemantik}

\subsection{Technischer Erfolg und Autorisierung}

Die zentrale Grenze ist
\begin{equation}
  \boxed{\mathrm{TRANSPORT\_ACK}\neq\mathrm{EFFECT\_ACK}}.
\end{equation}
Im Kandidaten gilt
\begin{equation}
 \operatorname{mayMutateExternally}(r,a)
 :=\operatorname{technicalSuccess}(r)
 \land\operatorname{externallyAuthorized}(a).
\end{equation}
Damit kann ein Exitcode \(0\) oder ein bestandener Test ohne eigene Autorität
keine Außenwirkung freigeben.

Die umfassendere QIK-VRT-Policy kennt fünf Zustände:
\begin{center}
\begin{tabular}{>{\ttfamily}l p{0.60\textwidth}}
\toprule
EFFECT\_NACK & Kein wirkungsprüfbarer Empfang. \\
EFFECT\_ACK\_CONTINUE & Prüfung darf fortgesetzt werden; keine gewöhnliche
Freigabe. \\
EFFECT\_ACK\_DONE & Alle für den exakten Scope geforderten Bedingungen sind
gebunden; einziger gewöhnlicher Freigabezustand. \\
EFFECT\_ACK\_ISOLATE & Kandidat für kontrollierte Prüfung separieren. \\
EFFECT\_ACK\_BLOCK & Kandidatenwirkung nicht fortsetzen. \\
\bottomrule
\end{tabular}
\end{center}

Diese LaTeX-Datei, ihr PDF und die Begleitartefakte bleiben
\code{EFFECT\_ACK\_CONTINUE}. Ihre Erzeugung ist keine GitHub-, Zenodo-,
IETF- oder sonstige externe Publikationsautorisierung.

\section{Repositorygebundener Belegstand}

Der herangezogene lokale Checkout war bei Erstellung dieser Fassung wie folgt
gebunden:

\begin{center}
\begin{tabularx}{\textwidth}{>{\bfseries}l X}
\toprule
Repository & \code{Goldkelch/qik-vrt} \\
Ref & lokales \code{main} \\
Commit & \code{2d4bb16278da34f8ee91eecd4f92a85ac02488ce} \\
Tree & \code{8f2d4ae07c41638af54811e51f154133d6b293b2} \\
Worktree & tracked und untracked sauber \\
Handoff & scopegebunden \code{PASS} \\
Integrität & 3183 klassifizierte Einträge, 3174 immutable Digests geprüft \\
Cachevertrag & 14/14 Komponenten, 100\% \\
Runtime & \code{CONTINUE}; Lean 4.19.0 und Lake nicht verfügbar \\
\bottomrule
\end{tabularx}
\end{center}

Der vorhandene v2-Formalisierungsreport weist für das gebundene Manuskript
20/20 Definitionsumgebungen, 20/20 theoremartige Umgebungen, 42 starke
quellgebundene Lean-Bindungen und 6 bedingte Bindungen aus. Dieser Scope ist
nicht mit dem neuen \VRTCore{}-Kandidaten gleichzusetzen.

Die frühere Textfamilie berichtet einen endlichen 21-Satz-Kandidaten. Im
aktuellen Checkout wurde kein eigenständiges Receipt gefunden, das genau diese
Zahl an den aktuellen Main bindet. Der hier vorgelegte neue 21-Satz-Kandidat
ist deshalb eigenständig nummeriert und bleibt bis zur Kernel-Ausführung offen.

\section{Menschliche Konsequenzen}

\subsection{Die Maschine für das Wessen und Warum}

Eine verantwortbare Erkenntnismaschine soll nicht nur Treffer oder plausible
Antworten liefern. Sie soll prüfen können:

\begin{quote}
Wessen Information ist das? Woher stammt sie? In welchem Kontext entstand
sie? Wie wurde sie verändert? Welche Unsicherheit bleibt? Warum soll sie
wirken? Wem nutzt sie? Wem kann sie schaden? Wer trägt Verantwortung? Welche
Rechte sind betroffen? Darf die Wirkung freigegeben werden?
\end{quote}

Diese Architektur ist kein Anspruch auf Allwissen. Sie ordnet Provenienz,
Kontext, Widerspruch, Unsicherheit, Wirkung, Verantwortung und Rechte so, dass
eine menschliche Entscheidung besser informiert und später auditierbar werden
kann.

\subsection{Das Recht auf einen Haltepunkt}

Ein System, das \code{OPEN}, \code{CONTINUE}, \code{ISOLATE} oder
\code{BLOCK} begründen kann, ist nicht schwächer als ein System, das immer
antwortet. Es besitzt eine zusätzliche verantwortliche Fähigkeit: Es kann den
eigenen Erkenntnis- oder Autorisierungsbereich erkennen.

Die normative Konsequenz lautet:
\begin{quote}
Menschen sollten nicht nur eine Antwort erhalten, sondern auch deren Herkunft,
Status, Unsicherheit, Widersprüche und Wirkungsgrenze erkennen können.
\end{quote}
Dies ist eine Sollensforderung, kein automatisch empirisch bewiesener Nutzen.

\subsection{Freiheit als Anschlussfähigkeit}

Die Formulierung
\begin{quote}
\textbf{Freiheit ist reale Anschlussfähigkeit im Hier und Jetzt.}
\end{quote}
ist eine interpretative QIK-VRT-These. Sie deutet Freiheit weder als
Kausallosigkeit noch als bloßes Epiphänomen, sondern als reale Möglichkeit,
zwischen unterscheidbaren und verantwortbaren Anschlüssen zu wählen, Gründe zu
prüfen und Korrektur wirksam zu machen. Ihr Status bleibt interpretativ, bis
philosophische, neurowissenschaftliche und empirische Brücken weiter
ausgearbeitet sind.

\subsection{Begründeter Stolz}

Die Arbeit darf ausdrücklich sagen, dass hier eine großartige Leistung
vorliegt. Der Grund ist nicht eine behauptete endgültige Wahrheit. Der Grund
ist die ungewöhnliche Verbindung aus:

\begin{itemize}
  \item einer weitreichenden kausal- und informationstheoretischen Intuition,
  \item einer typisierten Erkenntnisordnung,
  \item einer maschinenlesbaren Claimsprache,
  \item formalen Beweiskernen und negativen Übertreibungstests,
  \item Provenienz- und Korrekturspur,
  \item physikalischen Nichtaussagen und
  \item einer verantwortungsgebundenen Wirkungsgrenze.
\end{itemize}

\begin{quote}
\textbf{Aus einer großen Idee ist eine prüfbare, korrigierbare und teilweise
kernelgeprüfte Forschungsarchitektur geworden. Das ist eine großartige
Leistung. Darauf darf man stolz sein.}
\end{quote}

\section{Falsifikations- und Forschungsprogramm}

\subsection{Formale nächste Schritte}

\begin{enumerate}
  \item Die exakten Bytes des Lean-Kandidaten mit Lean 4.19.0 ausführen.
  \item Axiome, \code{sorry}/\code{admit} und Toolchain dynamisch auditieren.
  \item Ein Receipt an Quell-SHA-256, Werkzeugversion und alle 21 Deklarationen
  binden.
  \item Einen Parser für die EBNF implementieren und positive sowie negative
  Korpora prüfen.
  \item Die AST-Semantik auf die Lean-Datentypen abbilden und
  Statusaufwertungen negativ testen.
\end{enumerate}

\subsection{Mathematisch-physikalische nächste Schritte}

\begin{enumerate}
  \item Eine genaue Kategorie oder Ordnungsstruktur für Anschlussfähigkeit
  festlegen.
  \item Notwendige und hinreichende Bedingungen für einen stabilen
  Minkowski-Grenzfall formulieren.
  \item Existenz, Nicht-Eindeutigkeit und Stabilität der Rekonstruktion prüfen.
  \item Erst danach den Übergang zu allgemeinen lorentzschen Strukturen
  untersuchen.
  \item Eine quantitative Korrespondenz zu Prozessmatrizen, Quantum Switch oder
  Causal-Set-Dynamiken angeben.
  \item Messbare Abweichungen von bestehenden Modellen und klare
  Falsifikationsprotokolle definieren.
\end{enumerate}

\subsection{Widerlegungsbedingungen}

Die zentrale Forschungsthese wäre in ihrem jeweiligen Scope geschwächt oder
widerlegt, wenn beispielsweise gezeigt wird, dass:

\begin{itemize}
  \item die behauptete relationale Struktur keine wohldefinierte oder
  nichttriviale Semantik besitzt,
  \item jede vorgeschlagene physikalische Brücke lediglich bekannte Modelle
  umbenennt, ohne neue Erklärung oder Vorhersage,
  \item der klassische Grenzfall unter den angegebenen Voraussetzungen nicht
  existiert oder instabil ist,
  \item die resultierende Struktur etablierte Experimente verletzt oder
  \item die Claim- und Wirkungstypisierung ihre eigenen Statusgrenzen nicht
  zuverlässig erhalten kann.
\end{itemize}

\section{Explizite Nichtaussagen}

Diese Fassung behauptet nicht:

\begin{itemize}
  \item eine fertige Quanten- oder Quantengravitationstheorie,
  \item eine physische Identität zwischen Repositorygraph und Raumzeit,
  \item vollzogene Minkowski- oder allgemeine Lorentz-Emergenz,
  \item physische Retrokausalität oder operative Rückwärtssignalisierung,
  \item dass Prozessmatrizen oder Quantum Switch QIK-VRT empirisch bestätigen,
  \item dass die 21 neuen Sätze bereits kernelgeprüft sind,
  \item dass formale Konsistenz empirische Wahrheit erzeugt,
  \item unabhängiges Peer Review oder Feldkonsens,
  \item repositoryweites \code{PASS}, \code{FINAL\_PASS} oder
  \code{EFFECT\_ACK\_DONE} oder
  \item eine externe Publikation, einen Commit, Merge, Release oder eine
  Zenodo-Mutation.
\end{itemize}

\section{Schluss}

Die wissenschaftliche Pointe ist nicht, Sequenz abzuschaffen. Sie besteht
darin, Kausalität nicht länger aus Sequenz zu erschleichen. Die menschliche
Pointe ist nicht, Maschinen zu allwissenden Autoritäten zu machen. Sie besteht
darin, Herkunft, Unsicherheit, Widerspruch, Verantwortung und Wirkungsgrenze
zu Teilen ihrer Architektur zu machen.

\begin{quote}
\textbf{Unterschied macht Information möglich.}\\
\textbf{Messung bindet Information an Kontext.}\\
\textbf{Wirkung wird durch Relation erklärbar.}\\
\textbf{Kausalität ist Relation, nicht Sequenz.}\\
\textbf{Raumzeit ist ein zu rekonstruierender klassischer Grenzfall.}\\
\textbf{Verantwortung beginnt dort, wo Wirkung möglich wird.}
\end{quote}

Die daraus entstandene Forschungsarchitektur ist eine großartige Leistung.
Ihre wissenschaftliche Stärke zeigt sich darin, dass sie nicht nur ihre
Beweise, sondern auch ihre offenen Grenzen maschinenlesbar macht.

\begin{flushright}
\textbf{q.e.d.}\\
\textbf{Ingolf Lohmann}
\end{flushright}

\clearpage
\appendix

\section{Vollständiger Lean-Kandidat}

Der folgende Quelltext wird aus der Begleitdatei eingebunden. Seine Aufnahme in
das PDF ist kein Kernel-Receipt.

\lstinputlisting[language=Lean]{VRTCore_RelationalCausality_Candidate.lean}

\clearpage
\section{Vollständige EBNF-Grammatik}

\lstinputlisting[language={}]{VRTCore_Syntax.ebnf}

\section{Beleg- und Statusmatrix}

Die maschinenlesbare Datei \code{VRTCore\_CLAIM\_MATRIX.json} bindet alle 21
Kandidatensätze, die wissenschaftlichen Primäranschlüsse, die menschlichen
Claims und die offenen Schließbedingungen. Ein Kandidatensatz trägt dort den
aktuellen Typ \code{OPEN}, den Zieltyp \code{FORMAL\_PROVED} und den Status
\code{FORMAL\_CANDIDATE\_UNVERIFIED\_IN\_THIS\_RUNTIME}.

\section{Primärliteratur}

\begin{thebibliography}{99}

\bibitem{Oreshkov2012}
O.~Oreshkov, F.~Costa und Č.~Brukner,
``Quantum correlations with no causal order,''
\emph{Nature Communications} 3, 1092 (2012).
\href{https://doi.org/10.1038/ncomms2076}{doi:10.1038/ncomms2076}.

\bibitem{Chiribella2013}
G.~Chiribella, G.~M.~D'Ariano, P.~Perinotti und B.~Valiron,
``Quantum computations without definite causal structure,''
\emph{Physical Review A} 88, 022318 (2013).
\href{https://doi.org/10.1103/PhysRevA.88.022318}{doi:10.1103/PhysRevA.88.022318}.

\bibitem{Araujo2015}
M.~Araújo et al.,
``Witnessing causal nonseparability,''
\emph{New Journal of Physics} 17, 102001 (2015).
\href{https://doi.org/10.1088/1367-2630/17/10/102001}{doi:10.1088/1367-2630/17/10/102001}.

\bibitem{Goswami2018}
K.~Goswami et al.,
``Indefinite Causal Order in a Quantum Switch,''
\emph{Physical Review Letters} 121, 090503 (2018).
\href{https://doi.org/10.1103/PhysRevLett.121.090503}{doi:10.1103/PhysRevLett.121.090503}.

\bibitem{Bombelli1987}
L.~Bombelli, J.~Lee, D.~Meyer und R.~D.~Sorkin,
``Space-time as a causal set,''
\emph{Physical Review Letters} 59, 521--524 (1987).
\href{https://doi.org/10.1103/PhysRevLett.59.521}{doi:10.1103/PhysRevLett.59.521}.

\bibitem{Malament1977}
D.~B.~Malament,
``The class of continuous timelike curves determines the topology of
spacetime,''
\emph{Journal of Mathematical Physics} 18, 1399--1404 (1977).
\href{https://doi.org/10.1063/1.523436}{doi:10.1063/1.523436}.

\bibitem{PurvesShort2019}
T.~Purves und A.~J.~Short,
``Nonclassically causal correlations without backwards-in-time signaling,''
\emph{Physical Review A} 99, 022101 (2019).
\href{https://doi.org/10.1103/PhysRevA.99.022101}{doi:10.1103/PhysRevA.99.022101}.

\bibitem{VanDerLugt2023}
T.~van der Lugt, J.~Barrett und G.~Chiribella,
``Device-independent certification of indefinite causal order in the quantum
switch,''
\emph{Nature Communications} 14, 5811 (2023).
\href{https://doi.org/10.1038/s41467-023-40162-8}{doi:10.1038/s41467-023-40162-8}.

\end{thebibliography}

\clearpage
\vspace*{\fill}
\begin{center}
\begin{minipage}{0.78\textwidth}
\begin{center}
{\Large\bfseries Kausalität ist Relation, nicht Sequenz.}\par
\vspace{0.45em}
{\large Verantwortung beginnt dort, wo Wirkung möglich wird.}
\end{center}
\vspace{0.9em}\hrule\vspace{0.9em}
{\small\noindent\textbf{Quellen- und Wirkungsgrenze.}
Der ursprüngliche Artikel und die frühere QIK-VRT-Textfamilie bleiben
unverändert. Diese Datei ist eine neue, getrennte Fassung. Copyright
\textcopyright{} 2026 Ingolf Lohmann. Eine externe Veröffentlichung,
Rechteänderung oder technische Außenwirkung ist mit ihrer Erstellung nicht
verbunden.\par}
\vspace{0.9em}\hrule\vspace{0.9em}
\begin{center}
\code{EFFECT\_ACK\_CONTINUE}\\[0.8em]
\textbf{q.e.d. · Ingolf Lohmann}
\end{center}
\end{minipage}
\end{center}
\vspace*{\fill}

\end{document}
